% s-296-p296-m001; original PDF 296, 297, 298
\Aphorism{s-296-p296-m001}{355}
Wer sich eingesteht, wie transcendent für seine Zeichendeuterkunst die Geheimschrift des Lebens ist, wie Gesundheit und Schönheit auf wunderlichen Schleichwegen gehen dürfen und müssen, wie oft auf Spinnfäden und Falterflügeln sich der Samen der köstlichsten Dinge rettet, wo überall er unvermuthet sein Erdreich findet und Wurzeln schlägt — wer diesem so fest und zart verknüpften Causalgewebe der Welt nachspürt, der wird immer mehr davon zurück-, darüber hinauskommen, aus solch einer klingenden Harmonie irgend ein Einzelnes hinwegzuwünschen, sei es für sich betrachtet noch so werthlos, missgeschaffen und hässlich. Gab es nicht, giebt es nicht auch für dies Einzelne eine Möglichkeit der Einfügung, irrt nicht irgendwo im Weltraum die complementäre Hälfte umher, mit der es einst, vor der Zersplitterung ins Irdische, ein Ganzes bildete und vielleicht einmal wieder bilden wird — genug! damit ist es schon gerechtfertigt gegen Götter und Teufel. Freilich sind auch diese Götter und Teufel gerechtfertigt, die sich an dem missgestalteten Einzelnen vergreifen; darfst du nichts aus dem Totum der Weltvorgänge forthaben wollen, so auch nicht den Willen in dir, der sich gegen diese allumfassende Toleranz zur Wehre setzt. Was ist das? bei der normativen Ausdeutung von Erkenntnissen scheint ein \Sperr{Kreis}, nämlich der \Foreign{circulus vitiosus,} unvermeidlich zu sein; der Willensimpuls, den ein Theorem ausstrahlt, wird durch \Sperr{iterative} Anwendung desselben Theorems wieder zerstört. Wir haben keinen Grund, hier mit Beispielen zu sparen. Was ist Fatalismus? Die als Willensmotiv wirkende Erkenntniss von der Nothwendigkeit alles Geschehens. Alles geht, wie es gehen muss, sagt der Türke; \Sperr{folglich} ist es gleichgültig, ob und wie ich handle. Nun aber weiter: es ist auch gleichgültig, ob du dieses „folglich“ anerkennst oder nicht; magst du dich vom Fatum normativ bestimmen lassen oder nicht, beidemale erfüllst du selber nur das Fatum. Indem du ein Verhalten gegenüber der erkannten Causalität suchst, rückst du unaufhaltsam im causalen Uhrgang aller Dinge mit fort. \Sperr{Folglich} erlaubt dir der Fatalismus zweiter Ordnung, den Fatalismus erster Ordnung zu ignoriren und so zu handeln, als ob alles Geschehen in deiner Macht stünde, und auch von diesem neuen „folglich“ brauchst du nicht Kenntniss zu nehmen und kannst einen Fatalismus dritter Ordnung, der nur um einen Kreisprocess vom Fatalismus erster Ordnung abweicht, in deinem Verhalten zum Ausdruck bringen. — In der schönen grünen Zeit, als der \Sperr{Determinismus} noch etwas Neues war und als gefährlich-revolutionäre Errungenschaft modernen Denkens galt, haben schwärmerische Köpfe auf Grund der Unfreiheit des menschlichen Willens das \Sperr{Recht zu strafen} geleugnet. Auch hier liegt ein geschlossener Kreis vor, nur dass diese Strafrechtsleugner im ersten Quadranten stehen blieben; kann der Verbrecher nichts dafür, wenn er sündigt, so kann die Gesellschaft nichts dafür, wenn sie ihn straft — das Recht zu strafen besteht vor dem Determinismus ebensowenig wie das Recht, dem Strafenden in den Arm zu fallen. — Einen ähnlichen Cirkel hat der Christ mit seiner „allgemeinen Sündhaftigkeit“ erfunden. Das Bewusstsein einer guten That ist sündhaft, das Sichfreifühlen von Sünde ist ärgste Verstocktheit des Sünders; Reue und Zerknirschung ist der dem Menschen angemessene Normalzustand. Also ein chronischer Gewissensdruck, den man wegen seiner Gleichmässigkeit nicht mehr empfindet. Dem Christen wird das beständige Gefühl zugemuthet, vom rechten Wege abgewichen zu sein: sobald er auf dem rechten Wege zu sein glaubt, ist er sicher auf dem falschen. Er mag mithin thun und lassen, was er will, es ist immer dasselbe, nämlich immer Sünde — wenn alle Wege gleich falsch sind, sind auch alle gleich richtig. Dieses Christenthum ist die Vorschule zum sittlichen Nihilismus; der Christ macht aus seiner Moral, wie B. Constant sagt, \Foreign{une masse compacte et indivisible pour qu’elle se mêle le moins possible avec ses actions et le laisse libre dans tous les détails.} — Die genannten Einzelfälle, denen sich als neuestes Beispiel Nietzsches „Ethik der ewigen Wiederkunft“ anreihen könnte, geben zu verstehen, was sich fast von selbst versteht, dass sich aus Erkenntnissen von irgendwelcher \Sperr{Allgemeinheit} keine Normen des moralischen Verhaltens ablesen lassen.
\par

Allgemeinheit ist Nivellement, Einebnung der Unebenheiten, Ausgleichung der Gegensätze, Aufhebung der specifischen, rangordnenden, daher allein willenanregenden Besonderheiten der Dinge; ihr ethisches Correlat ist immer eine Form des \Sperr{Indifferentismus}, nie eine Form der Motivation.
\par

